\documentclass{article}
\usepackage[UTF8]{ctex}
\usepackage{amsmath,amssymb,amsfonts}
\usepackage{enumitem}
\usepackage{geometry}
\geometry{a4paper, margin=1in}
% === 错题本排版优化（不参与题目标识） ===
\setlength{\parindent}{0pt}
\setlength{\parskip}{0.35em plus 0.10em minus 0.08em}
\setlength{\emergencystretch}{3em}
\setlist{itemsep=0.25em, topsep=0.40em, parsep=0pt}
\allowdisplaybreaks[2]
% === 错题本排版优化结束 ===
\begin{document}
\section*{湘潭大学2025年数学分析真题}
\section*{197．湘潭大学2025年数学分析试卷}

\subsection{题目1}
1． \(\lim _{n \rightarrow \infty}\left(1-\frac{1}{2^{2}}\right)\left(1-\frac{1}{3^{2}}\right) \cdots\left(1-\frac{1}{n^{2}}\right)\) ．

\subsection{题目2}
2． \(\lim _{x \rightarrow \infty} x^{\frac{3}{2}}(\sqrt{x+1}+\sqrt{x-1}-2 \sqrt{x})\) ．

\subsection{题目3}
1． \(\int_{0}^{100 \pi} \sqrt{1-\cos 2 x} \mathrm{~d} x\) ．

\subsection{题目4}
2． \(\lim _{n \rightarrow \infty} \int_{0}^{1} \frac{x^{n}}{1+x} \mathrm{~d} x\) ．

三．（15 分）设 \(\left\{y_{n}\right\}\) 是严格单调增加的正无穷大量，且 \(\lim _{n \rightarrow \infty} \frac{x_{n}-x_{n-1}}{y_{n}-y_{n-1}}=0\) ，则 \(\lim _{n \rightarrow \infty} \frac{x_{n}}{y_{n}}=0\) ．

四．（10 分）证明不等式：\(x-\frac{x^{2}}{2}<\ln (1+x)<x, x>0\) ．

五．（15 分）讨论反常积分的敛散性： \(\int_{0}^{+\infty} \frac{\ln (1+x)}{x^{p}} \mathrm{~d} x, p \in \mathbb{R}\) ．

六．（10 分）计算重积分 \(\iint_{D} x y^{2} \mathrm{~d} x \mathrm{~d} y\) ，其中 \(D\) 为抛物线 \(y^{2}=2 p x\) 和直线 \(x=\frac{p}{2}(p>0)\) 所围的区域．

七．（15 分）设 \(f(x)\) 在区间 \(I\) 上为下凸函数，证明：\(f(x)\) 在区间 \(I\) 的任一闭子区间上有界．

八．（15 分）设函数 \(\varphi(x)\) 在 \([0,+\infty)\) 上存在二阶连续的导数，若 \(\lim _{x \rightarrow+\infty} \varphi(x)\) 存在，且 \(\varphi^{\prime \prime}(x)\) 有界．试证： \(\lim _{x \rightarrow+\infty} \varphi^{\prime}(x)=0\) ．

九．（15 分）设 \(f(x)=\sum_{n=0}^{\infty} \frac{1}{2^{n}+x}\) ．\\
（1）（9 分）证明：\(f(x)\) 在 \([0,+\infty)\) 上可导，且一致连续．\\
（2）（6 分）证明：反常积分 \(\int_{0}^{+\infty} f(x) \mathrm{d} x\) 发散．

附加题。（15 分）计算曲面积分

\[
\iint_{S}\left(y-x^{3}\right) \mathrm{d} y \mathrm{~d} z+\left(y^{3}-z\right) \mathrm{d} z \mathrm{~d} x+2 \mathrm{~d} x \mathrm{~d} y
\]

其中 \(S\) 是雉面 \(z^{2}=x^{2}+y^{2}(0 \leq z \leq 2)\) ，方向取下侧。
\end{document}
